\documentclass[12pt,a4paper]{article}

% ── Packages ──────────────────────────────────────────────────────────────
\usepackage[margin=1in]{geometry}
\usepackage{graphicx}
\usepackage{amsmath,amssymb}
\usepackage{booktabs}
\usepackage{longtable}
\usepackage{array}
\usepackage{multirow}
\usepackage{caption}
\usepackage{subcaption}
\usepackage{float}
\usepackage{hyperref}
\usepackage{listings}
\usepackage{xcolor}
\usepackage{enumitem}
\usepackage{fancyhdr}
\usepackage{tikz}
\usetikzlibrary{shapes,arrows,positioning,fit,calc}
\usepackage{titlesec}

% ── Code listing style ───────────────────────────────────────────────────
\lstdefinestyle{verilog}{
    language=Verilog,
    basicstyle=\ttfamily\footnotesize,
    keywordstyle=\color{blue}\bfseries,
    commentstyle=\color{green!50!black}\itshape,
    stringstyle=\color{red},
    numbers=left,
    numberstyle=\tiny\color{gray},
    numbersep=5pt,
    frame=single,
    breaklines=true,
    tabsize=4,
    showstringspaces=false,
    captionpos=b
}

% ── Header / Footer ──────────────────────────────────────────────────────
\pagestyle{fancy}
\fancyhf{}
\fancyhead[L]{Team 14 -- Pipelined RISC-V Processor}
\fancyhead[R]{IPA Project Phase II Report}
\fancyfoot[C]{\thepage}
\renewcommand{\headrulewidth}{0.4pt}

% ── Title ─────────────────────────────────────────────────────────────────
\title{%
    \vspace{-1cm}
    \textbf{Introduction to Processor Architecture}\\[6pt]
    \Large Pipelined RISC-V Processor -- Phase II Report\\[4pt]
    \large Team 14
}
\author{
    Member 1 \and Member 2 \and Member 3 \and Member 4
}
\date{\today}

% ══════════════════════════════════════════════════════════════════════════
\begin{document}
\maketitle
\thispagestyle{empty}
\newpage

\tableofcontents
\newpage

% ══════════════════════════════════════════════════════════════════════════
%  1.  OVERVIEW
% ══════════════════════════════════════════════════════════════════════════
\section{Overview}

This report covers the additions made on top of the sequential RISC-V processor to convert it into a 5-stage pipelined processor.
The base datapath blocks (ALU, Control Unit, Register File, Data Memory, Instruction Memory, Immediate Generator, ALU Control, Adder, ShiftLeft1, MUX) remain functionally identical to the sequential design.

\subsection{What Changed from the Sequential Design}
The following were added or modified:
\begin{enumerate}
    \item \textbf{Four sets of pipeline registers} between the five stages (IF/ID, ID/EX, EX/MEM, MEM/WB).
    \item \textbf{Forwarding Unit} -- resolves data hazards by bypassing results from later stages back to the EX stage ALU inputs.
    \item \textbf{Hazard Detection Unit} -- detects load-use and branch-related hazards that cannot be resolved by forwarding alone; stalls the pipeline by inserting bubbles.
    \item \textbf{Branch resolution moved to ID stage} -- a dedicated \texttt{equalityComparator} resolves \texttt{beq} in the Decode stage (rather than the EX stage), reducing the branch penalty from 2 flush cycles to 1.
    \item \textbf{Branch Forwarding Unit} -- provides forwarding paths to the branch comparator so that the ID-stage comparison uses up-to-date values.
    \item \textbf{Store-data forwarding (bonus)} -- a special forwarding path for store instructions whose data register depends on a value being written back simultaneously.
    \item \textbf{4-to-1 MUX} (\texttt{mux4x2.v}) -- new multiplexer module used for forwarding selection (4 instances).
    \item \textbf{PC enable signal} -- the PC module now has an \texttt{enable} input (\texttt{PCWrite}) that can freeze the PC during stalls.
\end{enumerate}

\subsection{Naming Convention for Pipeline Signals}
All pipeline register signals follow a prefix convention that makes them traceable in the codebase:
\begin{itemize}
    \item \textbf{\texttt{IFID*}} -- internal signals of the IF/ID pipeline register (e.g., \texttt{IFIDIR}, \texttt{IFIDPC}).
    \item \textbf{\texttt{IDEX*}} -- internal signals of the ID/EX register (e.g., \texttt{IDEXrs1}, \texttt{IDEXALUOp}, \texttt{IDEXrd}).
    \item \textbf{\texttt{EXMEM*}} -- internal signals of the EX/MEM register (e.g., \texttt{EXMEMALUResult}, \texttt{EXMEMrd}).
    \item \textbf{\texttt{MEMWB*}} -- internal signals of the MEM/WB register (e.g., \texttt{MEMWBReadData}, \texttt{MEMWBrd}).
    \item \textbf{\texttt{BMIDEX*}} -- ``Bubble Mux'' outputs that feed into ID/EX. When a bubble is asserted, these are zeroed to insert a NOP (e.g., \texttt{BMIDEXRegWrite}).
    \item Control signals \texttt{Forward*} come from forwarding units; hazard signals \texttt{PCWrite}, \texttt{IFIDEnable}, \texttt{Bubble} come from the Hazard Detection Unit.
\end{itemize}

% ══════════════════════════════════════════════════════════════════════════
%  2.  PIPELINE REGISTERS
% ══════════════════════════════════════════════════════════════════════════
\section{Pipeline Registers}

All four pipeline registers are implemented as \texttt{reg} declarations inside \texttt{processor.v} and are updated on \texttt{posedge clk}. On \texttt{reset}, instruction registers (\texttt{IFIDIR}, \texttt{IDEXIR}) are initialized to the NOP encoding \texttt{0x00000013} (\texttt{addi x0, x0, 0}), and all write-enable and memory-enable control signals are cleared to prevent false writes.

\subsection{IF/ID Register}
\begin{center}
\begin{tabular}{lcp{7cm}}
    \toprule
    \textbf{Signal} & \textbf{Width} & \textbf{Description} \\
    \midrule
    \texttt{IFIDIR} & 32 & Fetched instruction. Used inside the ID stage as \texttt{Instruction} via a continuous assign. \\
    \texttt{IFIDPC} & 64 & PC of the fetched instruction (used for branch target computation). \\
    \bottomrule
\end{tabular}
\end{center}

\textbf{Special behavior:}
\begin{itemize}
    \item When \texttt{IFIDEnable = 0} (stall), both fields hold their previous values.
    \item When \texttt{IFFlush = 1} (taken branch), \texttt{IFIDIR} is overwritten with the NOP instruction to nullify the wrongly fetched instruction.
\end{itemize}

\subsection{ID/EX Register}
\begin{center}
\begin{tabular}{lcp{7cm}}
    \toprule
    \textbf{Signal} & \textbf{Width} & \textbf{Source (from ID stage)} \\
    \midrule
    \texttt{IDEXALUOp} & 2 & ALU operation specifying signal \\
    \texttt{IDEXALUSrc} & 1 & select signal for ALU input B \\
    \texttt{IDEXMemRead} & 1 & read enable for Memory  \\
    \texttt{IDEXMemWrite} & 1 & write enable for Memory \\
    \texttt{IDEXMemtoReg} & 1 & Source register address for Write-back \\
    \texttt{IDEXRegWrite} & 1 & Register write enable \\
    \midrule
    \texttt{IDEXrs1} & 64 & Read data from register \texttt{rs1} \\
    \texttt{IDEXrs2} & 64 & Read data from register \texttt{rs2} \\
    \texttt{IDEXImm} & 64 & Sign-extended immediate \\
    \texttt{IDEXIR} & 32 & Full instruction (for funct3/funct7) \\
    \texttt{IDEXrd} & 5 & Destination register address \\
    \texttt{IDEXrs1Addr} & 5 & Source register 1 address \\
    \texttt{IDEXrs2Addr} & 5 & Source register 2 address \\
    \bottomrule
\end{tabular}
\end{center}

\textbf{Note:} The six control signals pass through the \textbf{Bubble Mux} before being latched.
When \texttt{Bubble = 1}, all six are zeroed, effectively converting the instruction in ID/EX into a NOP.
The data fields (\texttt{IDEXrs1}, \texttt{IDEXrs2}, \texttt{IDEXImm}, etc.) are always written regardless of bubble, since a NOP with zero control signals has no side effects.

\subsection{EX/MEM Register}
\begin{center}
\begin{tabular}{lcp{7cm}}
    \toprule
    \textbf{Signal} & \textbf{Width} & \textbf{Source (from EX stage)} \\
    \midrule
    \texttt{EXMEMALUResult} & 64 & ALU output \\
    \texttt{EXMEMrs2} & 64 & Forwarded rs2 value (for store data), sourced from \texttt{ALUm2Op}, (it's the latest value of rs2 for EX stage, so we use it) \\
    \texttt{EXMEMrs2Addr} & 5 & rs2 register address (for SD forwarding) \\
    \texttt{EXMEMrd} & 5 & Destination register \\
    \texttt{EXMEMMemRead} & 1 & Memory read enable \\
    \texttt{EXMEMMemWrite} & 1 & Memory write enable \\
    \texttt{EXMEMRegWrite} & 1 & Register write enable \\
    \texttt{EXMEMMemtoReg} & 1 & Write-back source \\
    \bottomrule
\end{tabular}
\end{center}

\textbf{Note:} \texttt{EXMEMrs2} stores the \emph{forwarded} rs2 value (\texttt{ALUm2Op}, output of the forwarding mux) rather than the raw \texttt{IDEXrs2}, so that forwarded data reaches the Data Memory correctly.

\subsection{MEM/WB Register}
\begin{center}
\begin{tabular}{lcp{7cm}}
    \toprule
    \textbf{Signal} & \textbf{Width} & \textbf{Source (from MEM stage)} \\
    \midrule
    \texttt{MEMWBReadData} & 64 & Data memory read output \\
    \texttt{MEMWBALUResult} & 64 & ALU result passed through \\
    \texttt{MEMWBrd} & 5 & Destination register \\
    \texttt{MEMWBRegWrite} & 1 & Register write enable \\
    \texttt{MEMWBMemtoReg} & 1 & Write-back source (0 = ALU, 1 = memory) \\
    \bottomrule
\end{tabular}
\end{center}

% ══════════════════════════════════════════════════════════════════════════
%  3.  DATA FORWARDING
% ══════════════════════════════════════════════════════════════════════════
\section{Data Forwarding}

\subsection{Forwarding Unit (\texttt{ForwardingUnit.v})}

The Forwarding Unit detects when an instruction in the EX stage needs a result that is still in the EX/MEM or MEM/WB pipeline register. It produces three signals:

\begin{itemize}
    \item \texttt{ForwardA} [1:0] -- selects ALU operand A source.
    \item \texttt{ForwardB} [1:0] -- selects ALU operand B source.
    \item \texttt{Forward\_sd} [1 bit] -- selects store-data source.
\end{itemize}

\textbf{Inputs:} \texttt{EXMEMRegWrite}, \texttt{MEMWBRegWrite}, \texttt{EXMEMrd}, \texttt{MEMWBrd}, \texttt{EXMEMrs2Addr}, \texttt{IDEXrs1Addr}, \texttt{IDEXrs2Addr}.

\subsubsection{Forwarding Selection}
\begin{center}
% Adjust the 8cm to fit your preference; 8cm usually fits well on A4 with 1in margins
\begin{tabular}{ccp{8cm}} 
    \toprule
    \textbf{\texttt{ForwardA/B}} & \textbf{MUX input} & \textbf{Condition (A uses rs1Addr, B uses rs2Addr)} \\
    \midrule
    \texttt{2'b00} & \texttt{IDEXrs1/2} & No hazard \\
    \addlinespace
    \texttt{2'b10} & \texttt{EXMEMALU} & EX hazard: \texttt{EXMEMRegWrite \&\& EXMEMrd $\neq$ 0 \&\& EXMEMrd == src} \\
    \addlinespace
    \texttt{2'b01} & \texttt{write\_back} & MEM hazard: \texttt{MEMWBRegWrite \&\& MEMWBrd $\neq$ 0 \&\& MEMWBrd == src} \\
    \bottomrule
\end{tabular}
\end{center}
EX hazard has higher priority over MEM hazard (checked first via \texttt{if}-\texttt{else}).

The forwarded values are routed through two \texttt{mux4x2} instances (\texttt{ALUm1} for operand A, \texttt{ALUm2} for operand B) before entering the ALU.

\subsubsection{Store-Data Forwarding (Bonus)}

When a store instruction reaches the MEM stage and its data register (\texttt{EXMEMrs2Addr}) matches the destination of the instruction currently in WB, the store data must be forwarded:

\texttt{Forward\_sd = 1} when \texttt{MEMWBRegWrite \&\& MEMWBrd $\neq$ 0 \&\& MEMWBrd == EXMEMrs2Addr}.

A \texttt{mux} is used to (\texttt{sdMUX}) select between \texttt{EXMEMrs2} (normal) and \texttt{write\_back} (forwarded) based on \texttt{Forward\_sd}. The output (\texttt{sdMUXOp}) feeds into the Data Memory's write port.

\subsection{Branch Forwarding Unit (\texttt{BranchForwardingUnit.v})}

Since the branch condition is resolved in the \textbf{ID stage} (not EX), and the branch comparator needs the latest register values, a separate forwarding unit is needed.

\textbf{Inputs:} \texttt{EXMEMRegWrite}, \texttt{MEMWBRegWrite}, \texttt{EXMEMrd}, \texttt{MEMWBrd}, \texttt{IFIDrs1} (= \texttt{Instruction[19:15]}), \texttt{IFIDrs2} (= \texttt{Instruction[24:20]}).

\textbf{Outputs:} \texttt{ForwardCompA} [1:0], \texttt{ForwardCompB} [1:0] -- same encoding as \texttt{ForwardA/B}. Two \texttt{mux4x2} instances (\texttt{compMuxA}, \texttt{compMuxB}) feed the forwarded values to the \texttt{equalityComparator}.

% ══════════════════════════════════════════════════════════════════════════
%  4.  HAZARD DETECTION UNIT
% ══════════════════════════════════════════════════════════════════════════
\section{Hazard Detection Unit (\texttt{HazardDetectionUnit.v})}

The Hazard Detection Unit detects data and control hazards that \emph{cannot} be resolved by forwarding alone. When a stall is required, it asserts:
\begin{itemize}
    \item \texttt{PCWrite = 0} -- freezes the PC.
    \item \texttt{IFIDEnable = 0} -- freezes the IF/ID register.
    \item \texttt{Bubble = 1} -- zeroes the control signals entering ID/EX (via the Bubble Mux), inserting a NOP.
\end{itemize}

\textbf{Inputs:} \texttt{IDEXMemRead}, \texttt{IDEXRegWrite}, \texttt{EXMEMMemRead}, \texttt{EXMEMrd}, \texttt{Instruction} (current IF/ID instruction), \texttt{IDEXIR} (ID/EX instruction).

The HDU internally derives:
\begin{lstlisting}[style=verilog, numbers=none]
wire [4:0] IDEXrd = IDEXIR[11:7];
wire [4:0] rs1 = Instruction[19:15];
wire [4:0] rs2 = Instruction[24:20];
wire isBEQ = (Instruction[6:0] == 7'b1100011);
\end{lstlisting}

\subsection{Stall Cases}

\subsubsection{Case 1: Load-Use Hazard (non-branch)}
When a \texttt{ld} instruction is in ID/EX and the next instruction in IF/ID reads the loaded register, forwarding cannot help because the data is not available until the end of MEM.

\textbf{For I-type (ld, addi):} only \texttt{rs1} is a source register.\\
\texttt{Stall if: IDEXMemRead \&\& IDEXrd $\neq$ 0 \&\& IDEXrd == rs1}

\textbf{For R-type / sd:} both \texttt{rs1} and \texttt{rs2} are source registers.\\
\texttt{Stall if: IDEXMemRead \&\& IDEXrd $\neq$ 0 \&\& (IDEXrd == rs1 || IDEXrd == rs2)}

\subsubsection{Case 2: Branch N-1 Dependency}
When \texttt{beq} is in IF/ID and the \emph{immediately preceding} instruction (in ID/EX) writes to a register that \texttt{beq} reads. Because the branch comparator is in the ID stage, it needs the value one cycle before forwarding from EX/MEM is available.

\texttt{Stall if: isBEQ \&\& IDEXRegWrite \&\& IDEXrd $\neq$ 0 \&\& (IDEXrd == rs1 || IDEXrd == rs2)}

\begin{itemize}
    \item If the preceding instruction is an ALU instruction: 1 stall, then forward from EX/MEM.
    \item If the preceding instruction is a \texttt{ld}: 1st stall here; after the \texttt{ld} moves to EX/MEM, Case 3 triggers a 2nd stall.
\end{itemize}

\subsubsection{Case 3: Branch N-2 Load Dependency}
When \texttt{beq} is in IF/ID and a load instruction two positions earlier (now in EX/MEM) writes to a register that \texttt{beq} reads. The loaded data only becomes available at the end of MEM, so one more stall is needed before it can be forwarded from MEM/WB.

\texttt{Stall if: isBEQ \&\& EXMEMMemRead \&\& EXMEMrd $\neq$ 0 \&\& (EXMEMrd == rs1 || EXMEMrd == rs2)}

% ══════════════════════════════════════════════════════════════════════════
%  5.  CONTROL HAZARD (BRANCH) HANDLING
% ══════════════════════════════════════════════════════════════════════════
\section{Control Hazard Handling}

\subsection{Branch Resolution in the ID Stage (Bonus)}

Instead of resolving branches in the EX stage (which would require flushing 2 instructions), we resolve branches in the \textbf{ID stage} using a dedicated \texttt{equalityComparator}:

\begin{lstlisting}[style=verilog, numbers=none]
equalityComparator comp1(.a(compA), .b(compB), .equal(rs1EqualRs2));
\end{lstlisting}

The inputs \texttt{compA} and \texttt{compB} come from the Branch Forwarding MUXes (\texttt{compMuxA}, \texttt{compMuxB}), which provide up-to-date (possibly forwarded) register values.

The branch decision: \texttt{IFFlush = Branch \& rs1EqualRs2}.
\begin{itemize}
    \item When \texttt{IFFlush = 1}: the IF/ID register is loaded with NOP (\texttt{0x00000013}), and the PC is redirected to the branch target (\texttt{pcPlusImm = IFIDPC + ShiftLeft1(immediate)}).
    \item Only 1 instruction in the IF stage needs to be flushed, giving a \textbf{1-cycle branch penalty} (vs.\ 2 cycles if resolved in EX).
\end{itemize}

\subsection{Branch Resolution in ID Stage}
To optimize performance, branch instructions are resolved in the Instruction Decode (ID) stage. A dedicated \texttt{equalityComparator} evaluates the branch condition immediately after the operands are fetched. This design choice reduces the branch penalty to a single clock cycle, as only the instruction currently in the IF stage must be flushed if the branch is taken. This implementation uses a \texttt{BranchForwardingUnit} to ensure the comparator has access to updated values from the EX and MEM stages.

% ══════════════════════════════════════════════════════════════════════════
%  6.  SPECIFIC TEST CASES
% ══════════════════════════════════════════════════════════════════════════
\section{Test Cases from the Instruction Program}

The test program (\texttt{instructions\_exp.txt}) was designed to exercise every hazard scenario. Below we describe each scenario with the relevant instruction sequence and how the pipeline handles it.

\subsection{EX-to-EX Forwarding (1-cycle-ago dependency)}
\begin{lstlisting}[style=verilog, numbers=none]
add  x5, x5, x5      // writes x5
sub  x6, x0, x5      // reads x5 (EX hazard: forward EXMEMALUResult)
\end{lstlisting}
\texttt{ForwardB = 2'b10}: \texttt{sub}'s operand B is forwarded from EX/MEM (the ALU result of \texttt{add}).

\subsection{MEM-to-EX Forwarding (2-cycle-ago dependency)}
\begin{lstlisting}[style=verilog, numbers=none]
add  x16, x15, x14   // writes x16
add  x17, x14, x15   // (no dependency on x16)
add  x18, x16, x17   // reads x16 (MEM hazard: forward from MEM/WB)
\end{lstlisting}
\texttt{ForwardA = 2'b01}: \texttt{x16}'s value is forwarded from MEM/WB.

\subsection{Load-Use Stall followed by Forwarding}
\begin{lstlisting}[style=verilog, numbers=none]
ld   x7, 17(x0)      // loads x7 from memory
add  x0, x0, x0      // NOP (if not present, HDU stalls 1 cycle)
sd   x7, 0(x8)       // uses x7
\end{lstlisting}
The NOP inserted by the programmer avoids the stall. Without the NOP, the HDU would detect \texttt{IDEXMemRead \&\& IDEXrd == rs2} and stall for 1 cycle.

A case that \emph{does} trigger the load-use stall:
\begin{lstlisting}[style=verilog, numbers=none]
ld   x24, 121(x11)   // loads x24
sub  x25, x0, x24    // reads x24 immediately -> 1-cycle stall
\end{lstlisting}
HDU detects: \texttt{IDEXMemRead \&\& IDEXrd(x24) == rs2(x24)} $\Rightarrow$ stall 1 cycle. After the stall, \texttt{x24}'s value is forwarded from EX/MEM.

\subsection{Branch with N-1 ALU Dependency (1 stall + EX/MEM forward)}
\begin{lstlisting}[style=verilog, numbers=none]
add  x13, x5, x5     // writes x13
beq  x13, x8, 8      // reads x13 (N-1: stall 1 cycle, then
                      //   BFU forwards EXMEMALUResult to compA)
addi x31, x0, -1     // flushed if branch taken
\end{lstlisting}
\begin{enumerate}
    \item HDU detects: \texttt{isBEQ \&\& IDEXRegWrite \&\& IDEXrd(x13) == rs1(x13)} $\Rightarrow$ stall 1 cycle.
    \item After the stall, \texttt{add x13} is in EX/MEM. BFU sets \texttt{ForwardCompA = 2'b10}.
    \item \texttt{compA = EXMEMALUResult} (x5+x5 = 0xFC), \texttt{compB = x8} (from register file = 0xFC).
    \item \texttt{rs1EqualRs2 = 1} $\Rightarrow$ branch taken, \texttt{addi x31} is flushed. x31 stays 0.
\end{enumerate}

\subsection{Branch with N-1 Load Dependency (2 stalls + MEM/WB forward)}
\begin{lstlisting}[style=verilog, numbers=none]
ld   x26, 1(x0)      // loads x26
beq  x10, x26, 8     // reads x26 (N-1 load: 2 stalls needed)
sub  x31, x0, x15    // flushed if branch taken
\end{lstlisting}
\begin{enumerate}
    \item \textbf{Stall 1 (Case 2):} HDU detects \texttt{isBEQ \&\& IDEXRegWrite \&\& IDEXrd(x26) == rs2(x26)} $\Rightarrow$ stall. The \texttt{ld} moves to EX/MEM.
    \item \textbf{Stall 2 (Case 3):} HDU detects \texttt{isBEQ \&\& EXMEMMemRead \&\& EXMEMrd(x26) == rs2(x26)} $\Rightarrow$ stall. The \texttt{ld} moves to MEM/WB.
    \item After 2 stalls, BFU sets \texttt{ForwardCompB = 2'b01} (from MEM/WB).
    \item \texttt{compB = write\_back = MEMWBReadData} (loaded value 0x11), \texttt{compA = x10} (from register file = 0x11).
    \item \texttt{rs1EqualRs2 = 1} $\Rightarrow$ branch taken, \texttt{sub x31} is flushed. x31 stays 0.
\end{enumerate}

\subsection{Store-Data Forwarding (Bonus)}
\begin{lstlisting}[style=verilog, numbers=none]
ld   x20, 0(x11)     // loads x20
sd   x20, 121(x11)   // stores x20 (load-use stall, then forwarding)
ld   x21, 121(x11)   // loads the stored value back
\end{lstlisting}
The \texttt{sd} instruction depends on \texttt{x20} which was just loaded. After the load-use stall and forwarding, the correct value of \texttt{x20} is written to memory. Additionally, the \texttt{Forward\_sd} path handles cases where the store data comes from MEM/WB.

\subsection{Internal Register File forwarding}
\begin{lstlisting}[style=verilog, numbers=none]
ld   x2, 1(x0)       // writes x2
add  x0, x0, x0      // NOP
add  x0, x0, x0      // NOP
add  x3, x2, x0      // reads x2 (3 cycles later, WB and ID same cycle)
\end{lstlisting}
The Register File module implements a combinational ``write-before-read'' bypass to manage dependencies between the Write-Back (WB) and Instruction Decode (ID) stages. When the destination register address of an instruction in the WB stage matches a source register address in the ID stage, the internal logic routes the \texttt{write\_data} bus directly to the \texttt{read\_data} output. 


\subsection{Branch Not Taken (No Flush)}
If \texttt{rs1 $\neq$ rs2}, the branch is not taken: \texttt{IFFlush = 0}, and the instruction in IF proceeds normally with no penalty.

% ══════════════════════════════════════════════════════════════════════════
%  7.  WAVEFORM PLACEHOLDERS
% ══════════════════════════════════════════════════════════════════════════
\section{Waveforms}

\begin{figure}[H]
    \centering
    \includegraphics[width=0.9\textwidth]{case1.png}
    \caption{At 80ns (cycle 9): ForwardB=10, EXMEMrd=x4, ALUresult=0xA,\\ A case of EX/MEM to EX forwarding}
    \label{fig:forwarding-waveform}
\end{figure}

\begin{figure}[H]
    \centering
    \includegraphics[width=0.9\textwidth]{case2.png}
    \caption{At 90ns (cycle 10): ForwardA=01, ForwardB=01, ALUresult=0xFFFFFFFFFFFFFFFB,\\ A case of MEM/WB to EX forwarding}
    \label{fig:forwarding-waveform}
\end{figure}

\begin{figure}[H]
    \centering
    \includegraphics[width=0.9\textwidth]{case3.png}
    \caption{At 50ns (cycle 6): ForwardA=01, ForwardB=10, ALUresult=0xF. A case of EX/MEM to EX and MEM/WB to EX forwarding}
    \label{fig:forwarding-waveform}
\end{figure}

\begin{figure}[H]
    \centering
    \includegraphics[width=0.85\textwidth]{case4.png}
    \caption{Waveform Analysis at 140\,ns (Cycle 15): Branch detection ($rs1EqualRs2=1$) triggers a pipeline flush ($IFFlush=1$). By 150\,ns (Cycle 16), a NOP (0x00000013) is inserted into the IF/ID register and the PC jumps to the target instruction.}
    \label{fig:forwarding-waveform}
\end{figure}

% ══════════════════════════════════════════════════════════════════════════
%  8.  RESULTS
% ══════════════════════════════════════════════════════════════════════════
\section{Results}

\subsection{Functional Correctness}
The processor was verified against the expected register file output. All 32 register values match the expected values, confirming correct handling of all forwarding, stalling, and flushing scenarios.

\begin{center}
\begin{tabular}{llll}
    \toprule
    \textbf{Register} & \textbf{Value} & \textbf{Register} & \textbf{Value} \\
    \midrule
    x0  & \texttt{0x0000000000000000} & x16 & \texttt{0x0000000000000276} \\
    x1  & \texttt{0xffffffffffffff04} & x17 & \texttt{0x0000000000000276} \\
    x2  & \texttt{0x0000000000000011} & x18 & \texttt{0x0000000000000762} \\
    x3  & \texttt{0x0000000000000011} & x19 & \texttt{0x0000000000000762} \\
    x4  & \texttt{0xffffffffffffff04} & x20 & \texttt{0xffffffffffffff82} \\
    x5  & \texttt{0x000000000000007e} & x21 & \texttt{0xffffffffffffff82} \\
    x6  & \texttt{0xffffffffffffff82} & x22 & \texttt{0x0000000000000011} \\
    x7  & \texttt{0xffffffffffffff82} & x23 & \texttt{0xffffffffffffff82} \\
    x8  & \texttt{0x00000000000000fc} & x24 & \texttt{0xffffffffffffff82} \\
    x9  & \texttt{0xffffffffffffff82} & x25 & \texttt{0x000000000000007e} \\
    x10 & \texttt{0x0000000000000011} & x26 & \texttt{0x0000000000000011} \\
    x11 & \texttt{0x0000000000000011} & x27 & \texttt{0x0000000000000011} \\
    x12 & \texttt{0xffffffffffffff82} & x28 & \texttt{0xffffffffffffff82} \\
    x13 & \texttt{0x00000000000000fc} & x29 & \texttt{0x0000000000000012} \\
    x14 & \texttt{0x00000000000001f8} & x30 & \texttt{0xffffffffffffff82} \\
    x15 & \texttt{0x000000000000007e} & x31 & \texttt{0x0000000000000000} \\
    \bottomrule
\end{tabular}
\end{center}

\subsection{Cycle Count}
Total clock cycles: \textbf{60}. This includes the pipeline fill latency, stall cycles for load-use and branch hazards, and flush cycles for taken branches.

\subsection{Hazard Scenarios Verified}
\begin{center}
\begin{tabular}{p{5cm}p{9cm}}
    \toprule
    \textbf{Scenario} & \textbf{Verified By} \\
    \midrule
    EX-to-EX forwarding & \texttt{add x5} $\to$ \texttt{sub x6}, \texttt{add x8} $\to$ \texttt{sd x7} \\
    MEM-to-EX forwarding & \texttt{add x16} $\to$ NOP $\to$ \texttt{add x18} \\
    Load-use stall & \texttt{ld x24} $\to$ \texttt{sub x25} \\
    Branch + ALU stall (N-1) & \texttt{add x13} $\to$ \texttt{beq x13, x8} \\
    Branch + load stall (N-1, 2 stalls) & \texttt{ld x26} $\to$ \texttt{beq x10, x26} \\
    Branch forwarding (EX/MEM) & \texttt{beq} after 1 stall, forward from EX/MEM \\
    Branch forwarding (MEM/WB) & \texttt{beq} after 2 stalls, forward from MEM/WB \\
    Store-data forwarding & \texttt{ld x20} $\to$ \texttt{sd x20} $\to$ \texttt{ld x21} \\
    Write-before-read bypass & \texttt{ld x2} $\to$ NOP $\to$ NOP $\to$ \texttt{add x3, x2, x0} \\
    Branch flush (taken) & Both \texttt{beq} correctly flush next instruction \\
    Branch no-flush (not taken) & Tested implicitly by correct x31 = 0 \\
    \bottomrule
\end{tabular}
\end{center}

% ══════════════════════════════════════════════════════════════════════════
%  9.  BONUS FEATURES SUMMARY
% ══════════════════════════════════════════════════════════════════════════
\section{BONUS sections}
Beyond basic hazard handling, the design includes specific optimizations to improve cycle efficiency:

\begin{itemize}
    \item \textbf{Store-Data Bypassing:} A specialized forwarding path (\texttt{Forward\_sd}) was implemented to handle dependencies where a store instruction's data is produced by an instruction in the Write-Back stage. This avoids a 1-cycle stall in load-store sequences.
    \item \textbf{Early Evaluation:} By resolving branches in the ID stage, the stall penalty is halved compared to EX-stage resolution. This involves integrated stalling logic that accounts for dependencies $N-1$ and $N-2$ level instructions.
    \item \textbf{Integrated Hazard Logic:} The Hazard Detection Unit applies the minimum necessary stall cycles for each specific data dependency, and implements logic for branch as well as non-branch related Hazards.
\end{itemize}

% ══════════════════════════════════════════════════════════════════════════
%  10. NEW FILES ADDED
% ══════════════════════════════════════════════════════════════════════════
\section{Files specific to Pipelined Implementation}

\begin{center}
\begin{tabular}{ll}
    \toprule
    \textbf{File} & \textbf{Purpose} \\
    \midrule
    \texttt{ForwardingUnit.v} & Data forwarding (EX/MEM hazard detection) \\
    \texttt{HazardDetectionUnit.v} & Load-use and branch stall detection \\
    \texttt{BranchForwardingUnit.v} & Branch comparator forwarding \\
    \texttt{equalityComparator.v} & ID-stage branch condition evaluation \\
    \texttt{mux4x2.v} & 4-to-1 MUX for forwarding paths \\
    \texttt{pipe\_tb.v} & System-level testbench \\
    \bottomrule
\end{tabular}
\end{center}

% ══════════════════════════════════════════════════════════════════════════

% ══════════════════════════════════════════════════════════════════════════
%  12. TEAM CONTRIBUTIONS
% ══════════════════════════════════════════════════════════════════════════
\section{Team Contributions}

\begin{center}
\begin{tabular}{ll}
    \toprule
    \textbf{Member} & \textbf{Contribution} \\
    \midrule
    Member 1 & [Fill in] \\
    Member 2 & [Fill in] \\
    Member 3 & [Fill in] \\
    Member 4 & [Fill in] \\
    \bottomrule
\end{tabular}
\end{center}

\end{document}
